\documentclass[landscape,a4paper,8pt]{article}
\usepackage[margin=1.1cm]{geometry}
\usepackage{multicol,amsmath,amssymb,booktabs,xcolor}
\setlength{\columnsep}{0.7cm}
\setlength{\parindent}{0pt}
\newcommand{\topic}[1]{\vspace{0.4em}\textcolor{blue!60!black}{\large\bfseries #1}\par\smallskip}
\newcommand{\key}[1]{\textbf{#1}: }
\title{Signals and Systems\\\large Course Summary Sheet}
\author{Your Name \quad Course code: EE---000}
\date{}
\begin{document}
\maketitle
\begin{multicols*}{3}
\topic{Core signals}
\key{Linearity} $T\{a x_1+b x_2\}=aT\{x_1\}+bT\{x_2\}$.\\
\key{Time invariance} A shift in $x(t)$ causes the same shift in $y(t)$.
\key{Convolution} $(x*h)(t)=\int_{-\infty}^{\infty}x(\tau)h(t-\tau)\,d\tau$.
\key{Impulse response} For an LTI system, $y=x*h$ and $h=T\{\delta\}$.
\topic{Transforms}
\[
X(j\omega)=\int x(t)e^{-j\omega t}dt,\qquad
x(t)=\frac1{2\pi}\int X(j\omega)e^{j\omega t}d\omega.
\]
\key{Convolution theorem} $x*h\leftrightarrow XH$.
\key{Differentiation} $dx/dt\leftrightarrow j\omega X(j\omega)$.
\key{Sampling} $f_s>2f_{\max}$ prevents aliasing for band-limited signals.
\topic{Filters}
\begin{tabular}{@{}lll@{}}\toprule
Filter & $H(j\omega)$ & Cutoff\\\midrule
RC low-pass & $1/(1+j\omega RC)$ & $1/RC$\\
RC high-pass & $j\omega RC/(1+j\omega RC)$ & $1/RC$\\
\bottomrule\end{tabular}
\topic{Problem checklist}
\begin{enumerate}
\item State assumptions and units.
\item Draw the system boundary and identify the input/output.
\item Choose time or frequency domain; do not mix conventions.
\item Check limiting cases and dimensions.
\item Interpret the result in one sentence.
\end{enumerate}
\topic{Exam workspace}
\vspace{7cm}
\topic{Common pitfalls}
Sign errors in Fourier exponents; forgetting initial conditions; treating a
finite-duration sample as band-limited; reporting a magnitude without units.
\end{multicols*}
\end{document}
